% Part: incompleteness
% Chapter: representability-in-q
% Section: c-representable

\documentclass[../../../include/open-logic-section]{subfiles}

\begin{document}

% SEGMENT OLP-0647-S01
\olfileid{inc}{req}{cre}
\olsection{$C$ இலுள்ள சார்புகள் $\Th{Q}$ இல் பிரதிநிதித்துவப்படுத்தத்தக்கவை}

$C$ இலுள்ள ஒவ்வொரு சார்பும் $\Th{Q}$ இல் பிரதிநிதித்துவப்படுத்தத்தக்கது
என்று நாம் காட்ட வேண்டும். இறுதியில், $C$ இலுள்ள ஒவ்வொரு $k$-இடச்
சார்பு $f(x_0,\dots,x_{k-1})$ க்கும் அதைப் பிரதிநிதித்துவப்படுத்தும்
வாய்பாடு $!A_f(x_0,\dots,x_{k-1},y)$ ஒன்றை எவ்வாறு ஒதுக்குவது என்று
காட்ட வேண்டும்.

% SEGMENT OLP-0647-S02
\begin{lem}
இயல் எண்கள் $n$, $m$ கொடுக்கப்பட்டிருக்க, $n \neq m$ எனில்
$Q \Proves \num n \neq \num m$.
\end{lem}

% SEGMENT OLP-0647-S03
\begin{proof}
ஒவ்வொரு $m$ க்கும், $n \neq m$ எனில்
$Q \Proves \eq/[\num n][\num m]$ என்று காட்ட $n$ மீது
தொகுத்தறிதலைப் பயன்படுத்துவோம்.

அடிப்படை நிலையில் $n = 0$. $m$ ஆனது $0$ க்குச் சமமல்ல எனில், ஏதோ
இயல் எண் $k$ க்கு $m = k + 1$. நம்மிடம்
$\lforall[x][\eq/[0][x']]$ என்று கூறும் அடிகோள் உள்ளது.
அளவையடை அடிகோளைப் பயன்படுத்தி $x$ க்குப் பதிலாக $\num k$ ஐ இட்டால்,
$\eq/[0][\num k']$ என்று முடிக்கலாம். ஆனால் $\num k'$ என்பது
$\num m$ தான்.

தொகுத்தறிதல் படியில், கூற்று $n$ க்கு உண்மை என்று எடுகொண்டு,
$n+1$ ஐக் கருதுவோம். $m$ தன்னிச்சையான இயல் எண்ணாக இருக்கட்டும்.
இரண்டு வாய்ப்புகள் உள்ளன: $m = 0$, அல்லது ஏதோ $k$ க்கு
$m = k+1$. முதல் நேர்வை மேலே கண்டவாறே கையாளலாம். இரண்டாம் நேர்வில்
$n+1 \neq k+1$ என்று கொள்க. அப்போது $n \neq k$.
$n$ க்கான தொகுத்தறிதல் எடுகோளால்
$\Th{Q} \Proves \eq/[\num n][\num k]$. நம்மிடம்
$\lforall[x][\lforall[y][\eq[x'][y'] \lif \eq[x][y]]]$
என்று கூறும் அடிகோள் உள்ளது. அளவையடை அடிகோளைப் பயன்படுத்தி,
$\eq[\num n'][\num k'] \lif \eq[\num n][\num
  k]$ ஐப் பெறுகிறோம். கூற்றுத் தருக்கத்தைப் பயன்படுத்தி,
$\Th{Q}$ இல்
$\eq/[\num n][\num k] \lif \eq/[\num n'][\num k']$
என்று முடிக்கலாம். விதிப்படி உட்படுத்தலால்
$\eq/[\num n'][\num k']$ ஐப் பெறுகிறோம். $\num k'$ என்பது
$\num m$ என்பதால், இதுவே வேண்டியது.
\end{proof}

% SEGMENT OLP-0647-S04
இந்தத் துணைத்தேற்றம் கூறுவது மிக அதிகமல்ல: வெவ்வேறு எண்ணுருக்கள்
வெவ்வேறு பொருட்களைக் குறிக்கின்றன என்பதை $\Th{Q}$ நிறுவ முடியும்
என்பதே அதன் சாரம். எடுத்துக்காட்டாக, $0'' \neq 0'''$ ஐ
$\Th{Q}$ நிறுவுகிறது. ஆனால் இது பொதுவாகப் பொருந்துகிறது என்று
காட்டுவதற்குச் சற்றுக் கவனம் தேவை. நாம் தொகுத்தறிதலைப் பயன்படுத்தினாலும்,
அது $\Th{Q}$ க்கு \emph{வெளியே} நடைபெறும் தொகுத்தறிதல் என்பதையும்
கவனிக்க.

சுழியம், அடுத்தெண், கூட்டல், பெருக்கல், சமத்துவத்தின் சிறப்பியல்புச்
சார்பு, வீழ்த்தல்கள் ஆகியவற்றை நம்மால் பிரதிநிதித்துவப்படுத்த முடியும்.
ஒவ்வொரு நேர்விலும் பொருத்தமான பிரதிநிதித்துவ வாய்பாடு முற்றிலும்
நேரடியானது; எடுத்துக்காட்டாக, சுழியம் பின்வரும் வாய்பாட்டால்
பிரதிநிதித்துவப்படுத்தப்படுகிறது:
\[
y = 0,
\]
அடுத்தெண் பின்வரும் வாய்பாட்டால் பிரதிநிதித்துவப்படுத்தப்படுகிறது:
\[
x_0' = y,
\]
கூட்டல் பின்வரும் வாய்பாட்டால் பிரதிநிதித்துவப்படுத்தப்படுகிறது:
\[
x_0 + x_1 = y.
\]
தொடர்புடைய கூற்றுகளை $\Th{Q}$ நிறுவ முடியும் என்று காட்டுவதில்தான்
பணி உள்ளது; எடுத்துக்காட்டாக, மேலுள்ள வாய்பாடு கூட்டலைப்
பிரதிநிதித்துவப்படுத்துகிறது என்பதற்கு, ஒவ்வோர் இயல் எண் சோடி
$m$,~$n$ க்கும் $\Th{Q}$ பின்வருவதை நிறுவுகிறது என்று காட்ட வேண்டும்:
\[
\eq[\num n + \num m][\num {n+m}]
\]
மேலும் பின்வருவதையும் நிறுவுகிறது:
\[
\lforall[y][(\eq[(\num n + \num m)][y] \lif \eq[y][\num {n+m}])].
\]

% SEGMENT OLP-0647-S05
சேர்ப்பு பற்றி என்ன? $h$ பின்வருமாறு வரையறுக்கப்படுகிறது என்று கொள்க:
\[
h(x_0,\dots,x_{l-1}) = f(g_0(x_0,\dots,x_{l-1}), \dots,
g_{k-1}(x_0,\dots,x_{l-1})).
\]
இங்கு $f,g_0,\dots,g_{k-1}$ ஆகிய சார்புகளை முறையே
பிரதிநிதித்துவப்படுத்தும் $!A_f, !A_{g_0}, \dots,
!A_{g_{k-1}}$ என்ற வாய்பாடுகளை ஏற்கெனவே கண்டுள்ளோம். அப்போது
$h$ ஐப் பிரதிநிதித்துவப்படுத்தும் $!A_h$ என்ற வாய்பாட்டை,
$!A_h(x_0,\dots,x_{l-1},y)$ பின்வருவதாக வரையறுப்பதன் மூலம்
அமைக்கலாம்:
% SOURCE CORRECTION TA-INC-041: keep all representing conjuncts in the quantifier scope.
\begin{multline*}
  \lexists[z_0\dots][\lexists[z_{k-1}][(!A_{g_0}(x_0,\dots,x_{l-1},z_0) \land
  \dots \land !A_{g_{k-1}}(x_0,\dots,x_{l-1},z_{k-1}) \land\\
  !A_f(z_0,\dots,z_{k-1},y))]]
\end{multline*}

% SEGMENT OLP-0647-S06
இறுதியாக வரம்பற்ற தேடலைக் கருதுவோம். $g(x,\vec z)$ ஒழுங்கானதாகவும்,
$!A_g(x,\vec z,y)$ என்ற வாய்பாட்டால் $\Th{Q}$ இல்
பிரதிநிதித்துவப்படுத்தத்தக்கதாகவும் இருக்கட்டும். $f$ ஆனது
$f(\vec z) = \mu x \; g(x,\vec z)$ என்பதால் வரையறுக்கப்படட்டும்.
$f$ ஐப் பிரதிநிதித்துவப்படுத்தும் $!A_f(\vec z,y)$ என்ற
வாய்பாட்டைக் கண்டறிய விரும்புகிறோம். இயல்பான தேர்வு இதுதான்:
\[
!A_f(\vec z,y) \ident !A_g(y,\vec z,0) \land \lforall[w][(w < y
\lif \lnot !A_g(w,\vec z,0))].
\]

இது செயல்படுகிறது என்பதைச் சில கூடுதல் துணைத்தேற்றங்களின் உதவியால்
காட்டலாம்; எடுத்துக்காட்டாக:

\begin{lem}
  ஒவ்வொரு மாறி $x$ க்கும் ஒவ்வோர் இயல் எண் $n$ க்கும்,
  $\Th{Q}$ ஆனது $\eq[(x'
  + \num n)][(x + \num n)']$ ஐ நிறுவுகிறது.
\end{lem}

இது, $\Th{Q}$ ஆனது
$\lforall[x][\lforall[y][(x' + y) = (x +y)']]$
ஐ நிறுவுகிறது என்று கூறுவதைவிட வலுக்குறைந்தது என்பதை மீண்டும்
குறிப்பிட வேண்டும்; அதை நிறுவ முடியும் என்ற வலுவான கூற்று பொய்.

% SEGMENT OLP-0647-S07
\begin{proof}
வழக்கம்போல $n$ மீது தொகுத்தறிதலால் நிறுவுவோம். அடிப்படை நிலையில்
$n = 0$; $\Th{Q}$ ஆனது
$\eq[(x'+0)][(x+0)']$ ஐ நிறுவுகிறது என்று காட்ட வேண்டும்.
ஆனால் நம்மிடம் பின்வருவன உள்ளன:
\begin{eqnarray*}
& & \eq[(x' + 0)][x'] \quad \text{அடிகோள் 4 இன்படி} \\
& & \eq[(x + 0)][x] \quad \text{அடிகோள் 4 இன்படி} \\
& & \eq[(x+0)'][x'] \quad \text{இரண்டாம் வரியின்படி} \\
& & \eq[(x' + 0)][(x + 0)'] \quad \text{முதல், மூன்றாம் வரிகளின்படி}
\end{eqnarray*}
தொகுத்தறிதல் படியில், $\Th{Q}$ இல்
$\eq[(x' + \num n)][(x + \num n)']$ ஐ
வருவித்துள்ளோம் என்று எடுகொள்கிறோம். $\num{n+1}$ என்பது
$\num n'$ என்பதால், $\Th{Q}$ ஆனது
$\eq[(x' + \num n')][(x + \num n')']$
ஐ நிறுவுகிறது என்று காட்ட வேண்டும். பின்வரும் சமத்துவத் தொடரை
$\Th{Q}$ இல் வருவிக்கலாம்:
\begin{eqnarray*}
(x' + \num n') & \eq & (x' + \num n)' \quad \text{அடிகோள் 5 இன்படி} \\
& \eq & (x + \num n')' \quad \text{தொகுத்தறிதல் எடுகோளின்படி}
\end{eqnarray*}
\end{proof}

% SEGMENT OLP-0647-S08
\begin{lem}
\begin{enumerate}
\item $\Th{Q}$ ஆனது $\lnot (x < \num 0)$ ஐ நிறுவுகிறது.
\item ஒவ்வோர் இயல் எண் $n$ க்கும், $\Th{Q}$ பின்வருவதை நிறுவுகிறது:
\[
x < \num {n+1} \lif (\eq[x][\num 0] \lor \dots \lor \eq[x][\num n]).
\]
\end{enumerate}
\end{lem}

% SEGMENT OLP-0647-S09
\begin{proof}
1 ஐயும் 2 இன் ஒரு பகுதியையும் முறைசாரா வகையில் நிறுவுவோம்
(அதாவது முறைசார் !!{derivation} ஐ அமைப்பதற்கான குறிப்புகளை மட்டுமே
தருவோம்).

1 ஆம் பகுதிக்கு, $<$ இன் வரையறையின்படி
$\Th{Q}$ இல்
$\lnot \lexists[y][\eq[(y'+x)][0]]$
ஐ நிறுவ வேண்டும்; இது முதல்தர தருக்கத்தின் அடிகோள்களையும் விதிகளையும்
பயன்படுத்தினால்
$\lforall[y][\eq/[(y' +
    x)][0]]$
என்பதற்குச் சமம். எண்ணம் இதுதான்: $\eq[(y'+x)][0]$
என்று கொள்க. $x$ ஆனது 0 எனில், $\eq[(y' + 0)][0]$
கிடைக்கும். ஆனால் $\Th{Q}$ இன் அடிகோள் 4 இன்படி
$\eq[(y' + 0)][y']$; அடிகோள் 2 இன்படி
$\eq/[y'][0]$; இது முரண்.
$x$ ஆனது $0$ அல்ல எனில், அடிகோள் 3 இன்படி
$\eq[x][z']$ ஆகும் $z$ ஒன்று உள்ளது. ஆனால் அப்போது
$\eq[(y' + z')][0]$ கிடைக்கும். அடிகோள்~5 இன்படி
$\eq[(y' + z)'][0]$ கிடைக்கிறது; இதுவும் அடிகோள்~2 உடன்
முரண்படுகிறது. இரு நேர்வுகளிலும்
$\lforall[y][\eq/[(y' +x)][0]]$ பெறப்படுகிறது.

2 ஆம் பகுதிக்கு, $n$ மீது தொகுத்தறிதலைப் பயன்படுத்துக.
$n =0$ என்ற அடிப்படை நிலையைக் கருதுவோம். அப்போது
$x < \num 1 \lif \eq[x][\num
  0]$ என்று காட்ட வேண்டும். $x < \num 1$ என்க.
$<$ க்கான வரையறுக்கும் அடிகோளின்படி
$\lexists[y][\eq[(y'+x)][0']]$ கிடைக்கும்.
அப்பண்புடைய $y$ ஒன்று இருப்பதாகக் கொள்க; எனவே
$\eq[y'+x][0']$.

$\eq[x][0]$ என்று காட்ட வேண்டும். அடிகோள் 3 இன்படி,
$x$ ஆனது 0 அல்ல எனில், அது ஏதோ $z$ க்கு $z'$ ஆகும்.
அப்போது $\eq[(y' + z')][0']$ கிடைக்கும்.
$\Th{Q}$ இன் அடிகோள்~5 இன்படி
$\eq[(y'+z)'][0']$ கிடைக்கும். அடிகோள் 1 இன்படி
$\eq[(y' +
    z)][0]$ கிடைக்கும். ஆனால், வரையறையின்படி இதற்கு
$z < 0$ என்று பொருள்; இது 1 ஆம் பகுதிக்கு முரணானது.
\end{proof}

% SEGMENT OLP-0647-S10
\begin{explain}
கணிக்கத்தக்க சார்புகளின் கணத்தை, $\Th{Q}$ இல்
பிரதிநிதித்துவப்படுத்தத்தக்க சார்புகளின் கணமாக வகைப்படுத்தலாம் என்று
காட்டியுள்ளோம். உண்மையில், நிறுவல் இன்னும் பொதுவானது.
பிரதிநிதித்துவப்படுத்தத்தக்கது என்பதன் வரையறையிலிருந்து,
$\Th{Q}$ ஐ விரிவாக்கும் எந்தக் கோட்பாட்டிலும்
(அல்லது $\Th{Q}$ ஐப் பொருள்படுத்தக்கூடிய கோட்பாட்டிலும்)
கணிக்கத்தக்க சார்புகளைப் பிரதிநிதித்துவப்படுத்தலாம் என்பதை
எளிதில் காணலாம். மறுபுறம், நிறுவல் என்ற கருத்து கணிக்கத்தக்கதாக
இருக்கும் எந்த !!{derivation} அமைப்பிலும் பிரதிநிதித்துவப்படுத்தத்தக்க
ஒவ்வொரு சார்பும் கணிக்கத்தக்கது. ஆகவே, எடுத்துக்காட்டாக,
கணிக்கத்தக்க சார்புகளின் கணம், பியானோ எண்கணிதத்தில் அல்லது
செர்மெலோ--ஃபிராங்கல் கணக் கோட்பாட்டில்கூட
பிரதிநிதித்துவப்படுத்தத்தக்க சார்புகளின் கணமாக வகைப்படுத்தப்படலாம்.
கோடல் குறிப்பிட்டதுபோல இது சற்றே வியப்பூட்டுகிறது.
நிறுவத்தன்மையைப் பொறுத்தவரை, எந்தக் கோட்பாட்டைக் கருதுகிறோம்
என்பதைப் பொறுத்துக் கேள்விகள் மிகவும் மாறுகின்றன என்பதைக் காண்போம்:
தோராயமாகச் சொன்னால், அடிகோள்கள் வலுப்பெறும்போது கூடுதலாக நிறுவலாம்.
ஆனால் பரந்த வகை அடிகோள் கோட்பாடுகளில், பிரதிநிதித்துவப்படுத்தத்தக்க
சார்புகள் துல்லியமாகக் கணிக்கத்தக்க சார்புகளே.
\end{explain}

\end{document}
